\frontmatter                    % Mandatory

% The order of the following components should be preserved.  The order
% listed here is the order currently required by FoGS.
\maketitle                      % Mandatory

\makeatletter
The following individuals certify that they have read, and recommend to the College of Graduate Studies for acceptance, a thesis/dissertation entitled:

\vspace{2ex}

%do NOT put curly braces after \ul to let soul expand by removing braces https://tex.stackexchange.com/questions/496833/package-soul-underline-problem-with-macro-text
\textsc{\ul\@title}



\vspace{1ex}

\noindent submitted by \underline{\sc \@author} in partial fulfilment of the requirements of the degree of \underline{\@degreetitle}
\makeatother

\newlength{\linespace}
\setlength{\linespace}{.75cm} %change .75cm to .5cm for smaller space between signatures or 1cm if less people have to sign
\vspace{2\linespace}\smaller
%%%%%%%%%%%%%%%%%%%%%%%%%%%%%%%%%%%%%%%%%%%%%%%%%%%%%%%%%%%%%%%%%%%%%%%%%%%%%%%%%%
% UPDATE THE FOLLOWING AS PER YOUR THESIS COMMITTEE

\noindent\underline{Alex S. Hill, Irving K. Barber Faculty of Science} 
\\
\textbf{Supervisor}

\vspace{\linespace}

\noindent\underline{John Braun, Irving K. Barber Faculty of Science} 
\\
\textbf{Supervisory Committee Member}

\vspace{\linespace}

\noindent\underline{Tim Robishaw, Irving K. Barber Faculty of Science} 
\\
\textbf{Supervisory Committee Member}

\vspace{\linespace}

\noindent\underline{Jonathan Holzman, Faculty of Applied Science} 
\\
\textbf{University Examiner}

\vspace{\linespace}

\noindent\underline{Tamara Freeman, Irving K. Barber Faculty of Science} 
\\
\textbf{External Examiner}

% \vspace{2\linespace}

% \noindent\textbf{Additional Committee Members include:}

% \vspace{\linespace}

% \noindent\underline{Jane Doe, Irving K. Barber Faculty of Science} 
% \\
% \textbf{Supervisory Committee Member}

% \vspace{\linespace}

% \noindent\underline{Jane Doe, Irving K. Barber Faculty of Science} 
% \\
% \textbf{Supervisory Committee Member}

%%%%%%%%%%%%%%%%%%%%%%%%%%%%%%%%%%%%%%%%%%%%%%%%%%%%%%%%%%%%%%%%%%%%%%%%%%%%%%%%%% END COMMITTEE PAGE
\normalsize

\newpage


\begin{abstract}                % Mandatory -  maximum 350 words
The Canadian Hydrogen Observatory and Radio-transient Detector (CHORD) is a 512-dish radio interferometer telescope under construction at the Dominion Radio Astrophysical Observatory (DRAO) near Penticton, British Columbia. This work estimates the capabilities of CHORD for detecting Radio Recombination Lines (RRLs) in the frequency range of 300-1500 MHz. This work analyzes the Radio Frequency Interference (RFI) present at DRAO and finds that 79 Hydrogen Alpha RRL transitions are likely to be detectable. This work derives the theoretical sensitivity of CHORD to diffuse sources and explores the expected signal-to-noise gains by stacking multiple RRLs. We propose site locations to add additional dishes as part of Extension CHORD, and find that adding 32 dishes across two locations can improve angular resolution by 44\% while decreasing sensitivity by 25\%. Finally, we take published results of the Survey of Ionized Gas in the Galaxy Made with the Arecibo Telescope (SIGGMA), and estimate CHORD's detection capabilities given those results. We find that CHORD should have the sensitivity to match each of SIGGMA's detections using four observing days and can improve upon each SIGGMA detection in both sensitivity and angular resolution after 9 observing days.


%This is a sample thesis based on the \texttt{ubcthesis.cls} template from Michael Forbes. The thesis includes the additional style file \texttt{ubcostyle.sty} in accordance to the official standards for the UBCO College of Graduate Studies. This sample thesis together with the style files and templates produces a document that is officially accepted by the UBCO College of Graduate Studies. 

%If you need a package, look into ubcostyle.sty to see if it is not already loaded there. 
%See the file README.txt for additional instructions to produce the bibliography, index, and glossary automatically.
\end{abstract}

\renewcommand{\abstractname}{Lay Summary}

\begin{abstract}
The Canadian Hydrogen Observatory and Radio-transient Detector (CHORD) is a new radio telescope under construction at the Dominion Radio Astrophysical Observatory (DRAO) near Penticton, British Columbia. One of the science goals for the telescope is to map the extremely faint radio waves originating from distant gas clouds in our Milky Way Galaxy. To accomplish this, the telescope must be extremely sensitive at specific radio frequencies. In this thesis we analyze which frequencies are usable by CHORD and which are corrupted by interference from terrestrial radio emitters. We calculate the theoretical sensitivity which we expect CHORD to achieve and provide data processing strategies that can improve sensitivity. We study the impact of adding additional antennas to CHORD and propose a design which optimizes scientific outcomes. Finally, we compare the expected performance of CHORD to previously published surveys and find that CHORD is likely to be more precise and more sensitive, likely allowing for new scientific discoveries.
\end{abstract}

\chapter{Preface}
This thesis was wholly written by Duncan Cameron-Steinke, who also performed all of the analysis presented in this thesis. This work was supervised by Dr. Alex Hill, with additional guidance from Dr. John Braun and Dr. Tim Robishaw. Nicholas Bruce assisted in the technical description of the RFI monitor and provided guidance for data processing of the RFI monitor data. Dr. Théo Michelot previously served as a committee member and provided continued guidance and support.

Generative AI in the form of Claude Code Opus 4.8 was used to accomplish specific, tightly scoped, coding tasks, such as generating plots and debugging. All generated code was carefully reviewed and validated to ensure correctness and accuracy. Generative AI in the form of Claude Code Opus 4.8 was used to edit versions of this thesis for grammar, clarity and to review accuracy of stated concepts and mathematical derivations. Generative AI was used as a resource while researching this thesis to provide background and answer questions.

Generative AI was not used to write any sections of this thesis and was not used to accomplish any of the analysis presented in the thesis. % no suggestions

\newpage
\phantomsection \label{tableofcontent}%set anchor at right location
\addcontentsline{toc}{chapter}{\contentsname}
\tableofcontents                % Mandatory: generate toc
\newpage 
\phantomsection \label{listoftab}%set anchor at right location
\addcontentsline{toc}{chapter}{\listtablename}
\listoftables                   % Mandatory if thesis has tables
\newpage
\phantomsection \label{listoffig}%set anchor at right location
\addcontentsline{toc}{chapter}{\listfigurename}
\listoffigures                  % Mandatory if thesis has figures


\chapter{Acknowledgements}      % Optional
Thank you to my supervisor Dr. Alex Hill, for being consistently available throughout my studies to answer questions and provide feedback. Also thank you Alex, for allowing me the flexibility to pursue my ever-changing research interests and always helping me along the way. \\

Thank you to all the members of the UBCO/UofC ISM group, Dr. Anna Ordog, Dr. Tom Landecker, Dr. Jo-Anne Brown, Dr. Jennifer West, Dr. Rebecca Booth, Ciara Chisholm, Art, Dr. Josh MacEachern, and Veyanna Hassani, for the many reviews and discussions.

Thank you to the CHORD spectral lines group, Dr. Kristin Spekkens, Akanksha Bij, Hans Hopkinson, and Josh Goodeve, for your contributions.


\chapter{Dedication} % Optional
Thank you Rebeka for your unwavering support.



% Any other unusual prefactory material should come here before the
% main body.

% Now regular page numbering begins.
\mainmatter
